\chapter{Fundamentação Teórica}
\noindent
\section{Criptografia Leve}

A disseminação de dispositivos de Internet das Coisas (Internet of Things - IoT), sistemas embarcados e sensores de baixo consumo impôs restrições operacionais que os algoritmos criptográficos convencionais não atendem. O Advanced Encryption Standard (AES) e as funções de hash da família SHA foram projetados para hardware com recursos computacionais amplos. Seu desempenho degrada-se em microcontroladores de 8 ou 16 bits, em circuitos RFID e em sensores alimentados por bateria, nos quais a área de silício, o consumo de energia e a memória disponível são limitados \cite{nist_sp800232_2025}.
A Criptografia Leve (Lightweight Cryptography - LWC) corresponde ao subcampo da criptografia simétrica dedicado a esquemas que preservam garantias de segurança comparáveis às de algoritmos convencionais sob restrições de área de chip, consumo energético, memória RAM/ROM e latência \cite{nist_sp800232_2025}.

As restrições típicas dos dispositivos-alvo da LWC incluem: área de circuito reduzida, memória RAM e ROM limitadas a algumas centenas ou poucos milhares de bytes, ausência de fonte de energia contínua, o que impõe baixo consumo médio e de pico, e processadores sem suporte a instruções vetorizadas ou aceleração criptográfica em hardware. Esse perfil é representativo de etiquetas RFID, implantes médicos, sensores de redes IoT e transponders veiculares \cite{nist_sp800232_2025}.
Sob tais condições, um esquema de criptografia autenticada com dados associados (Authenticated Encryption with Associated Data - AEAD) precisa entregar confidencialidade e autenticidade com um custo computacional muito inferior ao exigido pelo AES em modo Galois/Counter Mode (AES-GCM), mantendo um nível de segurança considerado adequado para a aplicação.

O National Institute of Standards and Technology (NIST) iniciou em 2015 o Lightweight Cryptography Standardization Process, processo público e multirrodada com o objetivo de selecionar algoritmos de criptografia leve adequados a ambientes restritos nos quais os padrões convencionais (AES-GCM, SHA-2, SHA-3) apresentam custo elevado \cite{nist_sp800232_2025}.
O processo seguiu o modelo já empregado pelo NIST em competições anteriores, com submissão pública de candidatos, rodadas sucessivas de eliminação e relatórios de status documentando os critérios de avaliação em cada etapa: o Status Report on the First Round (NIST IR 8268), o Status Report on the Second Round (NIST IR 8369) e o Status Report on the Final Round (NIST IR 8454) \cite{nist_sp800232_2025}. A avaliação considerou três eixos: segurança criptográfica, custo de implementação em hardware e software, e desempenho em cenários de uso representativos dos dispositivos-alvo.

Em fevereiro de 2023, o NIST anunciou a decisão de padronizar a família Ascon, projetada por Dobraunig, Eichlseder, Mendel e Schläffer, como resultado do processo de seleção \cite{nist_lwc_finalreport_2023}. A família Ascon havia sido originalmente submetida à competição CAESAR (Competition for Authenticated Encryption: Security, Applicability, and Robustness) em 2014, processo que selecionou Ascon-128 e Ascon-128a como recomendação primária para o caso de uso de criptografia autenticada leve em seu portfólio final, anunciado em 2019 \cite{cesar}. A formalização do padrão ocorreu com a publicação do NIST SP 800-232, cuja versão draft foi disponibilizada em novembro de 2024 para comentário público e cuja versão final foi publicada em agosto de 2025. O documento especifica quatro algoritmos: Ascon-AEAD128, esquema de criptografia autenticada baseado em nonce com força de segurança de 128 bits no cenário de chave única; Ascon-Hash256, função de hash que produz saída de 256 bits com força de segurança de 128 bits; e as funções de saída expansível (eXtendable Output Functions - XOF) Ascon-XOF128 e Ascon-CXOF128, esta última com suporte a string de personalização \cite{nist_sp800232_2025}.

O esquema Ascon-AEAD128 constitui um dos dois algoritmos avaliados nesta dissertação, descrito em detalhe na Seção 2.3. O segundo algoritmo, GIFT-COFB, não foi padronizado pelo NIST, mas permaneceu como finalista do mesmo processo de seleção, sendo descrito na Seção 2.4. A escolha desse par decorre diretamente do resultado da padronização: Ascon-AEAD128 corresponde ao algoritmo recomendado pelo NIST para uso em dispositivos restritos, enquanto GIFT-COFB representa um esquema de projeto distinto que alcançou nível equivalente de escrutínio criptográfico no mesmo processo. Essa equivalência de escrutínio é o que torna interpretável a comparação ciphertext-only conduzida nos Capítulos 4 e 5.



\section{Cifragem Autenticada com Dados Associados (AEAD)}

A criptografia simétrica clássica resolve a confidencialidade: dado um texto em claro e uma chave, produz um criptograma do qual um adversário sem a chave não consegue extrair informação sobre o conteúdo original. Essa garantia, isoladamente, não impede que o criptograma seja alterado em trânsito sem que o destinatário perceba. Um esquema de Cifragem Autenticada com Dados Associados (Authenticated Encryption with Associated Data - AEAD) combina as duas propriedades em um único primitivo: confidencialidade do texto em claro e autenticidade tanto do texto em claro quanto de um conjunto adicional de dados que acompanha a mensagem sem ser cifrado, os dados associados.

%Formalmente, um esquema AEAD define um algoritmo de cifragem que recebe uma chave secreta `K`, um nonce `N`, dados associados `A` de tamanho variável e um texto em claro `P` de tamanho variável, produzindo um criptograma `C`, com `|C| = |P|`, e uma tag de autenticação `T`:

%```
%AEAD.enc(K, N, A, P) = (C, T)
%```

O algoritmo de decifragem recebe a chave, o nonce, os dados associados, o criptograma e a tag, e retorna o texto em claro somente se a tag for válida; caso contrário, retorna uma indicação de falha, sem expor qualquer texto em claro parcial \cite{nist_sp800232_2025}. Os dados associados ilustram a diferença central de um esquema AEAD em relação à cifragem simples: cabeçalhos de protocolo, identificadores de sessão ou metadados de roteamento podem precisar circular em claro, para que intermediários da rede os leiam, mas qualquer alteração nesses campos deve invalidar a tag e ser detectada no destino \cite{nist_sp800232_2025}.

A autenticidade entregue por um esquema AEAD depende do uso correto do nonce. Um nonce deve ser distinto para cada operação de cifragem sob a mesma chave, condição que garante que textos em claro idênticos produzam criptogramas diferentes a cada execução \cite{nist_sp800232_2025}. A reutilização de um par (nonce, dados associados) sob a mesma chave compromete as garantias de confidencialidade e, em alguns esquemas, as de integridade, o que torna a geração e o controle do nonce um requisito operacional do esquema, e não um detalhe de implementação.

Ascon-AEAD128 e GIFT-COFB constituem os dois algoritmos avaliados nesta dissertação por atenderem à mesma definição de esquema AEAD descrita acima, com chave de 128 bits e força de segurança de 128 bits no cenário de chave única. Essa equivalência funcional é a condição que permite a comparação proposta no Capítulo 4: os dois algoritmos recebem entradas do mesmo tipo, sob as mesmas restrições de uso de nonce e chave, e produzem criptogramas com a mesma garantia teórica de indistinguibilidade frente a um adversário sem acesso à chave. Qualquer diferença estatística detectável entre os criptogramas dos dois esquemas, em um cenário ciphertext-only, não decorreria de diferenças na classe de problema que cada um resolve, mas de características residuais da construção interna de cada algoritmo, descritas nas Seções 2.3 e 2.4.


\section{Ascon-AEAD128}

Ascon-AEAD128 corresponde à instância de criptografia autenticada com dados associados (AEAD) da família Ascon, padronizada pelo NIST no SP 800-232 \cite{nist_sp800232_2025}. O algoritmo recebe uma chave secreta de 128 bits, um nonce de 128 bits, dados associados de tamanho variável e um texto em claro de tamanho variável, e produz um criptograma de mesmo tamanho do texto em claro e uma tag de autenticação de 128 bits, conforme a definição geral de AEAD apresentada na Seção 2.2 \cite{nist_sp800232_2025}.

A construção interna do algoritmo segue o modo esponja (sponge construction), no qual um estado interno é sucessivamente atualizado por uma função de permutação, absorvendo blocos de entrada e, ao final, liberando os blocos de saída \cite{nist_sp800232_2025}. O estado interno de Ascon-AEAD128 tem 320 bits, organizado em cinco palavras de 64 bits, e divide-se em uma porção de taxa (rate) de 128 bits, na qual ocorre a absorção dos dados associados e do texto em claro, e uma porção de capacidade (capacity) de 192 bits, que permanece oculta ao longo do processamento e sustenta a segurança do esquema \cite{nist_sp800232_2025}. A relação entre taxa e capacidade é o parâmetro central de qualquer construção esponja: quanto maior a capacidade em relação à taxa, maior a margem de segurança teórica contra ataques de distinção, ao custo de processar menos bits por chamada de permutação.

O núcleo da construção é a permutação Ascon-p, aplicada ao estado interno em duas variantes: Ascon-p[12], com 12 rodadas, usada nas fases de inicialização e finalização, e Ascon-p[8], com 8 rodadas, usada na absorção dos blocos de dados associados e de texto em claro \cite{nist_sp800232_2025}. Cada rodada da permutação segue a estrutura de Rede de Substituição-Permutação (Substitution-Permutation Network - SPN) e compõe três camadas em sequência: adição de constante, substituição não linear e difusão linear \cite{nist_sp800232_2025}. A camada de adição de constante insere, a cada rodada, um valor constante distinto em uma das cinco palavras do estado, o que rompe a simetria entre rodadas. A camada de substituição aplica, em paralelo, uma S-box de 5 bits a cada posição do estado, tomando um bit de cada uma das cinco palavras como entrada, o que introduz a não linearidade necessária para resistir a ataques lineares e diferenciais. A camada de difusão linear aplica a cada palavra do estado uma combinação de rotações de bits, espalhando a influência de cada bit de entrada sobre múltiplos bits de saída dentro da mesma palavra \cite{nist_sp800232_2025}.

O processamento de Ascon-AEAD128 ocorre em quatro fases sequenciais. Na inicialização, o estado é montado pela concatenação de um valor inicial fixo, da chave e do nonce, e atualizado por uma execução de Ascon-p[12], seguida pela combinação da chave com a porção final do estado. No processamento dos dados associados, cada bloco de 128 bits é combinado com a porção de taxa do estado e seguido por uma execução de Ascon-p[8], com um bit de separação de domínio aplicado ao final dessa fase para distinguir o processamento de dados associados do processamento subsequente do texto em claro. No processamento do texto em claro, cada bloco é combinado com a porção de taxa do estado para gerar o bloco correspondente do criptograma, com Ascon-p[8] aplicada entre blocos consecutivos. Na finalização, a chave é novamente combinada ao estado, uma execução de Ascon-p[12] é aplicada, e a tag de autenticação de 128 bits é extraída pela combinação da chave com a porção final do estado resultante \cite{nist_sp800232_2025}.

Ascon-AEAD128 compartilha a mesma permutação Ascon-p com as demais funções da família padronizada no SP 800-232: a função de hash Ascon-Hash256 e as funções de saída expansível Ascon-XOF128 e Ascon-CXOF128 empregam exclusivamente Ascon-p[12], em um modo esponja com taxa reduzida a 64 bits, ao passo que Ascon-AEAD128 emprega tanto Ascon-p[12] quanto Ascon-p[8], com taxa de 128 bits \cite{nist_sp800232_2025}. Essa partilha de permutação foi um dos critérios considerados pelo NIST na padronização: um único bloco criptográfico de baixo custo, reutilizado em múltiplas funcionalidades, reduz a área de implementação necessária em dispositivos restritos, ponto retomado na Seção 2.1.

A relevância de Ascon-AEAD128 para esta dissertação decorre diretamente de sua condição de padrão NIST: representa o algoritmo recomendado pelo processo de padronização de criptografia leve para o caso de uso de criptografia autenticada, e constitui o primeiro elemento do par avaliado nesta dissertação. A construção esponja descrita nesta seção contrasta com a construção de GIFT-COFB, descrita na Seção 2.4, baseada em uma cifra de bloco operando em modo de realimentação, o que situa os dois algoritmos em famílias de projeto distintas dentro do mesmo processo de seleção.


\section{GIFT-COFB}

GIFT-COFB constitui o segundo algoritmo avaliado nesta dissertação, projetado por Banik, Chakraborti, Iwata, Minematsu, Nandi, Peyrin, Sasaki, Sim e Todo, e submetido ao processo de padronização de criptografia leve do NIST descrito na Seção 2.1 \cite{Banik2021}. O algoritmo alcançou a rodada final do processo, ao lado de Ascon, sem ser selecionado como padrão, condição que define sua posição nesta dissertação: representante de um segundo eixo de projeto, com escrutínio criptográfico equivalente ao de Ascon-AEAD128, mas baseado em uma construção estrutural distinta \cite{nist_lwc_finalreport_2023}.

Ao contrário de Ascon-AEAD128, que opera sobre uma permutação em modo esponja, GIFT-COFB combina dois componentes: a cifra de bloco GIFT-128 e o modo de operação COFB (Combined Feedback). GIFT-128 é uma Rede de Substituição-Permutação (Substitution-Permutation Network - SPN) com bloco de 128 bits, chave de 128 bits e 40 rodadas, projetada como sucessora da cifra PRESENT com ganhos de eficiência em área de hardware e maior resistência a ataques lineares \cite{Banik2017}. Cada rodada de GIFT-128 compõe três etapas: SubCells, na qual uma S-box invertível de 4 bits é aplicada em paralelo a cada um dos 32 nibbles do estado; PermBits, permutação fixa de bits que redistribui a saída das S-boxes entre nibbles distintos, sem uso de portas lógicas adicionais; e AddRoundKey, na qual uma subchave derivada do escalonamento de chaves é combinada ao estado por XOR, junto com uma constante de rodada \cite{Banik2017}.

O modo COFB define como GIFT-128 processa dados associados e texto em claro de tamanho arbitrário a partir do bloco fixo de 128 bits da cifra subjacente. O modo é de passagem única, executando uma chamada à cifra de bloco por bloco de entrada, e dispensa a operação de decifragem da cifra de bloco, exigindo apenas a função de cifragem tanto para cifrar quanto para decifrar \cite{Banik2021}. Essa característica reduz a área de implementação em hardware, por eliminar a necessidade de circuitos para a transformação inversa da cifra de bloco. A cada bloco processado, a saída da chamada à cifra de bloco é combinada por uma função de retroalimentação (feedback function) com o bloco de entrada seguinte, mecanismo que dá nome ao modo e que encadeia o estado de uma chamada à próxima sem exigir um estado interno maior que o da própria cifra de bloco \cite{Banik2021}.

Na instanciação usada em GIFT-COFB, o nonce tem 128 bits e a tag de autenticação tem 128 bits, com a função de retroalimentação introduzindo uma perda de entropia de 1 bit em relação à versão original do modo COFB, ajuste que viabiliza um tamanho de nonce maior sob a mesma cifra de bloco de 128 bits \cite{Banik2021}. O esquema é classificado como inverse-free entre os finalistas do processo NIST que utilizam cifra de bloco como primitivo subjacente, condição compartilhada apenas com Romulus dentro desse grupo \cite{nist_lwc_finalreport_2023}.

A escolha de GIFT-COFB como segundo algoritmo do par avaliado nesta dissertação decorre dessa combinação de fatores: trata-se de um finalista do mesmo processo de padronização que produziu Ascon-AEAD128, construído sobre uma classe de primitivo distinta, cifra de bloco em modo de retroalimentação, em contraste com a construção em permutação e modo esponja descrita na Seção 2.3. Essa diferença de projeto interno, somada à equivalência de função AEAD e de nível de segurança estabelecida na Seção 2.2, define as condições sob as quais a comparação ciphertext-only dos Capítulos 4 e 5 é conduzida: qualquer padrão estatístico detectável entre os criptogramas dos dois esquemas decorreria de características residuais de suas construções internas, e não de uma assimetria na classe de problema que cada um resolve.



\section{Indistinguibilidade e o Cenário Ciphertext-Only}

A confidencialidade de um esquema criptográfico é normalmente formalizada por noções de indistinguibilidade: se a saída do algoritmo se comporta, para um adversário computacionalmente limitado, como uma sequência aleatória, então separar instâncias cifradas torna-se inviável. Shannon introduziu o ideal de segurança perfeita, no qual o criptograma, sem a chave, não revela nenhuma informação sobre o texto em claro \cite{shannon_1949}.
A criptografia computacional moderna relaxa esse ideal para o cenário de adversários eficientes: a base da hierarquia de definições é a indistinguibilidade sob ataque de texto-claro escolhido (IND-CPA), segundo a qual nenhum adversário probabilístico em tempo polinomial consegue distinguir, com vantagem não negligenciável, qual de duas mensagens de mesmo tamanho foi cifrada, mesmo podendo escolher as mensagens e observar seus criptogramas \cite{GoldwasserMicali1984PE}. Extensões mais fortes, IND-CCA1 e IND-CCA2, adicionam ao adversário acesso a um oráculo de decifragem antes e, no segundo caso, também depois do desafio, formando uma hierarquia em que IND-CCA2 implica IND-CCA1, que por sua vez implica IND-CPA.

A identificação algorítmica, problema central desta dissertação, decorre diretamente dessa formalização. Dado um conjunto de algoritmos $\mathcal{A} = \{A_1, \ldots, A_k\}$ e um conjunto de criptogramas $\mathcal{C}$ produzidos por um único $A_j \in \mathcal{A}$, a pergunta de pesquisa é se existe um classificador $f$ tal que $f(\mathcal{C}) = A_j$ com vantagem estatisticamente significativa sobre o acaso. Se dois algoritmos produzem saídas indistinguíveis de ruído sob o protocolo experimental adotado, não deveria existir um procedimento eficiente capaz de identificá-los de forma confiável a partir apenas das saídas. A atribuição algorítmica constitui, portanto, um caso particular do problema de distinção entre distribuições, e não um ataque que recupera chave, texto em claro ou qualquer informação além da origem do criptograma.

Modelos de ameaça descrevem as capacidades do adversário ao interagir com o sistema criptográfico. Os modelos clássicos incluem o ataque ciphertext-only (COA), no qual o adversário observa apenas criptogramas; o ataque de texto conhecido (KPA), no qual dispõe de pares texto-claro/criptograma; o ataque de texto escolhido (CPA), no qual escolhe os textos em claro a serem cifrados; e o ataque de texto-cifrado escolhido (CCA), no qual interage com um oráculo de decifragem. Esta dissertação adota o cenário ciphertext-only: os classificadores recebem exclusivamente os criptogramas, sem acesso à chave, ao nonce, ao texto em claro ou à tag de autenticação, e sem interação com oráculos de cifragem ou decifragem. Esse enquadramento delimita o problema de atribuição algorítmica como o caso mais restrito de distinção entre distribuições de saída, condição que torna qualquer resultado positivo de classificação particularmente informativo sobre vazamento residual da construção do algoritmo.

Em criptografia, afirmar que um criptograma "parece aleatório" significa, em termos operacionais, que suas amostras não exibem desvios estatisticamente detectáveis em relação às propriedades esperadas de uma sequência de bits com distribuição uniforme e independência aproximada. Essa noção empírica de aleatoriedade é distinta das definições formais de segurança apresentadas acima: um esquema pode produzir saídas que passam em diversos testes estatísticos e ainda ser inseguro por razões estruturais, e pode falhar em um teste específico sem que isso implique, de imediato, um ataque prático. Ainda assim, testes estatísticos, como os reunidos no NIST Statistical Test Suite \cite{Rukhin2010NIST80022}, oferecem um instrumento auxiliar para caracterizar empiricamente criptogramas e contextualizar padrões residuais. Múltiplos testes elevam a chance de falsos positivos, e o resultado depende da representação adotada (bit a bit, byte a byte, concatenação de mensagens, alinhamento de blocos) e do protocolo de amostragem; por isso, nesta dissertação, esses testes são tratados como ferramentas complementares, e não como substitutos de critérios formais de segurança.

A conexão entre essas noções teóricas e o uso de classificadores de aprendizado de máquina é direta: se existem desvios estatísticos sistemáticos na saída de um algoritmo, decorrentes de escolhas de projeto, modo de operação, empacotamento, padding, política de nonce ou artefatos de implementação, esses desvios podem atuar como sinal explorável por um classificador. O cenário ciphertext-only adotado nesta dissertação testa exatamente essa hipótese para o par Ascon-AEAD128 e GIFT-COFB, descritos nas Seções 2.3 e 2.4: ambos os esquemas, sob a teoria de segurança que fundamenta sua padronização e seleção como finalistas, deveriam produzir criptogramas indistinguíveis de ruído para um adversário sem acesso à chave, condição que os Capítulos 4 e 5 avaliam empiricamente.



\section{Aprendizado de Máquina Supervisionado para Classificação}

Um classificador de aprendizado de máquina supervisionado corresponde a uma função aprendida a partir de exemplos rotulados, que mapeia uma representação numérica de entrada x para uma classe y, produzindo uma decisão e, quando aplicável, uma probabilidade associada a essa decisão. No problema desta dissertação, x corresponde ao criptograma ou a características extraídas dele, e y corresponde ao algoritmo gerador, Ascon-AEAD128 ou GIFT-COFB. Esta seção apresenta, em nível conceitual, os classificadores empregados no Caminho A (Random Forest, SVM, LinearSVC e XGBoost) e as arquiteturas convolucionais empregadas como extratores de representação nos Caminhos B, C e D (CNN-1D e CNN-2D), descritos em detalhe no Capítulo 4.

Árvores de decisão particionam recursivamente o espaço de atributos em regiões homogêneas, com regras do tipo se-então sobre os valores de cada atributo. Uma árvore individual é interpretável e lida bem com mistura de tipos de variáveis, mas tende a alta variância: pequenas mudanças no conjunto de treinamento podem produzir árvores estruturalmente distintas. O Random Forest aborda essa limitação por agregação: treina um conjunto de árvores, cada uma sobre um subconjunto aleatório de amostras e de atributos, e combina as decisões individuais por votação, o que reduz a variância em relação a uma árvore isolada e melhora a capacidade de generalização sem exigir ajuste extenso de hiperparâmetros.

A Máquina de Vetores de Suporte (Support Vector Machine - SVM) busca o hiperplano que separa as classes maximizando a margem entre elas. Na versão linear, a fronteira de decisão é uma combinação linear dos atributos de entrada; com o kernel de Base Radial (Radial Basis Function - RBF), usado nesta dissertação, a SVM projeta implicitamente os dados em um espaço de maior dimensão, no qual fronteiras não lineares no espaço original tornam-se lineares, à custa de exigir ajuste cuidadoso dos hiperparâmetros de regularização (C) e de largura do kernel (γ).

Métodos de gradient boosting constroem o modelo de forma aditiva: cada árvore adicionada ao conjunto corrige, de forma iterativa, os erros residuais das árvores anteriores. O XGBoost implementa esse princípio com controles de regularização e otimizações de eficiência computacional, e é reconhecido por desempenho elevado em dados tabulares, característica que motiva seu uso nesta dissertação sobre o vetor de 29 descritores estatísticos extraídos de cada criptograma.

Random Forest, SVM e XGBoost operam sobre vetores de atributos calculados manualmente a partir do criptograma, descritos na Seção~\ref{sec:extracao-atributos}. Uma via alternativa de representação, empregada nos Caminhos B, C e D desta dissertação, dispensa o cálculo manual de descritores e aprende a representação diretamente a partir dos bytes do criptograma, por meio de Redes Neurais Convolucionais (Convolutional Neural Networks - CNN). Uma CNN aplica filtros convolucionais que percorrem a entrada e aprendem, durante o treinamento, quais padrões locais são relevantes para a tarefa, em vez de receber esses padrões como atributos pré-definidos pelo pesquisador.

A CNN unidimensional (CNN-1D) trata o criptograma como uma sequência de bytes e aplica filtros convolucionais ao longo dessa sequência, de forma análoga ao processamento de sinais temporais ou de áudio, capturando dependências locais entre bytes vizinhos e, por composição de camadas, padrões em escalas crescentes de comprimento. A CNN bidimensional (CNN-2D), por sua vez, opera sobre uma representação matricial do criptograma, neste caso um mapa de co-ocorrência de bigramas de bytes, e aplica filtros convolucionais sobre as duas dimensões dessa matriz, capturando padrões de coocorrência entre pares de valores de byte que uma representação puramente sequencial não expõe diretamente.

Em ambos os casos, a função das arquiteturas convolucionais nesta dissertação é estritamente a de extratoras de representação latente, e não de classificadoras fim-a-fim: a camada final de cada CNN produz um vetor de características aprendidas, extraído por uma função extract\_latent(), que é então fornecido a um classificador clássico --- Random Forest ou LinearSVC, isoladamente nos Caminhos B e C, e Random Forest ou XGBoost quando combinado às 307 características clássicas no Caminho D. Essa separação entre extração e classificação permite comparar, sob o mesmo protocolo experimental e os mesmos classificadores finais, se representações aprendidas diretamente dos bytes do criptograma carregam sinal que os descritores estatísticos manuais da Seção~\ref{sec:extracao-atributos} não capturam, sem confundir o ganho atribuível à representação com o ganho atribuível ao classificador.


\section{Seleção de Atributos}
\label{sec:selecao-atributos}

Um vetor de descritores com muitas dimensões e poucas amostras por classe favorece o sobreajuste e eleva o custo computacional dos classificadores, de modo que a redução de dimensionalidade antecede a etapa de treinamento em praticamente todo pipeline de aprendizado supervisionado sobre dados tabulares \cite{guyon2003}. Este trabalho recorre a três técnicas em sequência --- Informação Mútua, \emph{minimum Redundancy Maximum Relevance} (mRMR) e Boruta ---, complementadas por um índice de estabilidade baseado no coeficiente de Jaccard, descritos nesta seção.

A Informação Mútua (\emph{Mutual Information} --- MI) mede a dependência estatística entre uma variável aleatória $X$ (um descritor) e a variável-alvo $Y$ (a classe), sem assumir relação linear entre elas \cite{cover_thomas_2006}. Formalmente,
\begin{equation}
I(X;Y) = \sum_{x \in X}\sum_{y \in Y} p(x,y) \log\left(\frac{p(x,y)}{p(x)\,p(y)}\right),
\end{equation}
em que $p(x,y)$ é a distribuição conjunta e $p(x)$, $p(y)$ as distribuições marginais. $I(X;Y) = 0$ indica independência entre o descritor e a classe; valores maiores indicam maior dependência. Aplicada individualmente a cada descritor, a MI funciona como filtro univariado: ordena os candidatos por relevância isolada e descarta os que não carregam informação sobre o rótulo, sem levar em conta a relação entre os próprios descritores.

O filtro por MI, no entanto, não evita a seleção de descritores redundantes entre si: dois descritores altamente correlacionados podem receber pontuação de relevância semelhante e ser incluídos juntos, sem ganho de informação adicional. O mRMR corrige essa limitação ao combinar dois critérios no mesmo escore: a relevância de cada candidato frente à classe e a redundância dele frente aos descritores já selecionados \cite{peng_mrmr_2005}. Na formulação incremental, o próximo descritor $x_j$ a entrar no subconjunto $S$ é escolhido por
\begin{equation}
x_j = \argmax_{x_i \notin S} \left[ I(x_i; Y) - \frac{1}{|S|}\sum_{x_s \in S} I(x_i; x_s) \right],
\end{equation}
em que o primeiro termo recompensa a relevância e o segundo penaliza a redundância média em relação ao que já está em $S$. O procedimento é guloso e incremental: a cada passo, entra o candidato que melhor equilibra os dois critérios, até se atingir o número de descritores desejado.

Enquanto MI e mRMR produzem um subconjunto de tamanho fixo, o Boruta verifica, descritor a descritor, se a contribuição de cada um é estatisticamente distinguível do ruído \cite{kursa_boruta_2010}. O algoritmo duplica o conjunto original criando uma cópia \emph{sombra} de cada descritor, com os valores embaralhados entre as amostras, e treina uma Floresta Aleatória (\emph{Random Forest} --- RF) sobre o conjunto combinado (descritores reais e sombra). A importância de cada descritor real é então comparada, ao longo de várias iterações, com a importância máxima obtida entre os descritores sombra: um descritor real só é confirmado como relevante se superar essa referência de ruído de forma consistente, e é descartado se ficar sistematicamente abaixo dela. Por depender de reamostragem e da variabilidade intrínseca da RF, o conjunto de descritores confirmados pelo Boruta pode variar entre execuções e entre partições dos dados --- variação que este trabalho explora, no Capítulo~\ref{cap:resultados}, como indício adicional sobre a presença ou ausência de sinal real nos descritores avaliados.

Essa variação entre partições é quantificada pelo coeficiente de Jaccard, medida de similaridade entre dois conjuntos definida pela razão entre o tamanho da interseção e o tamanho da união:
\begin{equation}
J(A,B) = \frac{|A \cap B|}{|A \cup B|},
\end{equation}
em que $J(A,B) \in [0,1]$: $J=1$ indica conjuntos idênticos e $J=0$ indica conjuntos disjuntos. Aplicado aos subconjuntos de descritores selecionados em diferentes \emph{folds} de validação cruzada, o índice de Jaccard médio funciona como medida de estabilidade da seleção: quanto mais próximo de zero, menos os \emph{folds} concordam sobre quais descritores carregam sinal.

As quatro técnicas se combinam na etapa de seleção descrita na Seção~\ref{sec:metodologia-selecao}: MI e mRMR reduzem o vetor original a um subconjunto compacto dentro de cada \emph{fold} de treinamento, o Boruta verifica a estabilidade estatística desse subconjunto, e o coeficiente de Jaccard resume essa estabilidade em um único número, comparável entre os diferentes caminhos experimentais.

\section{Métricas e Validação Cruzada}

A avaliação de um classificador exige métricas de desempenho preditivo e um protocolo de validação que garanta independência entre dados de ajuste e de avaliação. Esta seção define as métricas e o esquema de validação cruzada empregados nesta dissertação, conceitos assumidos conhecidos a partir do Capítulo 4.

Em classificação binária, a acurácia mede a proporção de acertos sobre o conjunto avaliado. Sob classes balanceadas --- caso desta dissertação, com 30.000 amostras por algoritmo ---, acurácia e balanced accuracy coincidem; esta última, definida como a média das revocações por classe, generaliza a acurácia para cenários desbalanceados, nos quais a acurácia simples pode mascarar baixo desempenho na classe minoritária.

A medida F1 combina precisão (proporção de predições de uma classe que estão corretas) e recall (proporção de exemplos verdadeiros de uma classe corretamente identificados) pela média harmônica entre as duas:
\begin{equation}
F1 = 2 \times \frac{\text{precisão} \times \text{recall}}{\text{precisão} + \text{recall}}.
\end{equation}
Com mais de duas classes, ou quando se deseja peso igual entre elas, calcula-se o F1 por classe e toma-se a média aritmética, métrica denominada F1-macro. Nesta dissertação, binária e balanceada, o F1-macro funciona como métrica-alvo única; sob classificação ao acaso entre duas classes igualmente prováveis, seu valor esperado é 0,50, referência central na interpretação dos resultados.

A curva ROC (\emph{Receiver Operating Characteristic}) traça a taxa de verdadeiros positivos contra a de falsos positivos à medida que o limiar de decisão varia, e a área sob a curva (\emph{Area Under the Curve} --- AUC) resume esse traçado em um valor entre 0 e 1 \cite{fawcett2006}. AUC igual a 0,50 corresponde a desempenho ao acaso; AUC igual a 1,0, a separação perfeita entre as classes. Como o F1-macro, tem em 0,50 seu ponto de referência central, e só se aplica a classificadores que exportam probabilidade.

Quando um classificador reporta, além da classe prevista, a confiança associada a essa previsão, é possível avaliar se essa confiança corresponde à frequência real de acerto. O Erro de Calibração Esperado (\emph{Expected Calibration Error} --- ECE) quantifica esse descompasso: agrupam-se as predições em compartimentos por confiança reportada, calcula-se, em cada um, a diferença entre confiança média e acurácia observada, e combina-se essas diferenças ponderadas pelo tamanho de cada compartimento:
\begin{equation}
\text{ECE} = \sum_{m=1}^{M} \frac{|B_m|}{n} \left| \text{acc}(B_m) - \text{conf}(B_m) \right|,
\end{equation}
em que $B_m$ é o conjunto de predições no compartimento $m$, $|B_m|$ seu tamanho e $n$ o total de exemplos. ECE igual a zero indica calibração perfeita. No contexto desta dissertação, em que o resultado central é a ausência de sinal distintivo entre os dois algoritmos AEAD, o ECE serve como verificação complementar: um classificador sem informação útil deveria reportar baixa confiança, não confiança elevada associada a desempenho próximo do acaso.

Cada métrica reportada vem acompanhada de intervalo de confiança, obtido por reamostragem \emph{bootstrap} sobre o conjunto de teste \cite{efron1994}: reamostra-se, com reposição, um conjunto do mesmo tamanho do original, recalcula-se a métrica sobre essa reamostra, e repete-se o processo por $B=1.000$ iterações; o intervalo de 95\% corresponde aos percentis 2,5 e 97,5 da distribuição empírica resultante. O procedimento dispensa suposições paramétricas sobre a distribuição do erro.

O valor pontual de uma métrica não sustenta, por si só, que um classificador supera outro: a diferença observada pode decorrer de variação amostral. O teste de McNemar avalia essa hipótese comparando dois classificadores sobre o mesmo conjunto de teste, a partir da assimetria entre os exemplos em que os dois discordam \cite{mcnemar1947}; sob a hipótese nula de desempenho equivalente, a estatística resultante segue aproximadamente uma distribuição qui-quadrado com um grau de liberdade. Em comparações múltiplas sobre o mesmo conjunto de teste, o limiar de significância de cada teste é ajustado pela correção de Bonferroni, que divide o limiar nominal pelo número de comparações realizadas \cite{dunn1961}.

A validade dessas métricas depende da separação entre dados de ajuste e de avaliação. A validação cruzada particiona o conjunto em $k$ dobras, treina o modelo $k$ vezes --- cada vez com $k-1$ dobras para treinamento e a restante para validação --- e agrega o desempenho das $k$ rodadas. Essa forma padrão assume que cada exemplo pode ser distribuído independentemente entre dobras, suposição que falha quando exemplos compartilham origem comum cuja identidade não deveria estar disponível ao modelo: nesta dissertação, cada chave criptográfica gera 100 amostras por algoritmo, e um modelo que aprendesse padrões associados à chave, em vez de ao algoritmo, obteria desempenho artificialmente inflado caso exemplos da mesma chave aparecessem em treino e validação de uma mesma dobra.

O GroupKFold corrige isso ao particionar por grupo --- aqui, a chave criptográfica --- em vez de por exemplo individual, mantendo todos os exemplos de uma chave na mesma dobra e garantindo que nenhuma chave de validação tenha aparecido no treinamento daquela dobra. Esta dissertação aplica GroupKFold com 5 dobras sobre as 240 chaves reservadas para treinamento e validação, mantendo as 60 restantes em um conjunto de teste holdout, isolado de qualquer ajuste de modelo ou seleção de características (Capítulo 4). Essa combinação de key-holdout externo com GroupKFold interno é a condição metodológica que distingue o protocolo desta dissertação das avaliações da literatura discutidas no Capítulo 3, nas quais reaproveitamento de chave entre treino e teste produziu desempenho acima do que se sustenta sob separação rigorosa.


























